\documentclass{article}
\usepackage[utf8]{inputenc}
\usepackage[T1]{fontenc}
\usepackage{amsmath,amssymb}
\usepackage{graphicx}
\usepackage{booktabs}
\usepackage{hyperref}
\newcommand{\method}{V-JEPA~2}
\title{V-JEPA 2: Self-Supervised Video Models}
\author{A. Author \and B. Author}

\begin{document}
\maketitle

\begin{abstract}
We present \method{}, a self-supervised learning approach that learns from more than 1 million hours of internet video. % TODO: check number
The model reaches 77.3\% top-1 accuracy on Something-Something~v2 and achieves state-of-the-art results on action anticipation~\cite{damen2022}.
\end{abstract}

\section{Introduction}
\label{sec:intro}
A major challenge for modern AI is to learn to understand the world and act largely by observation~\cite{lecun2022}.
We train an encoder with a mask-denoising objective in representation space, i.e.\ we predict the embedding $s_y = f_\theta(y)$ of masked patches.
Our main \emph{contributions} are listed below.\footnote{Code is available at \url{https://github.com/facebookresearch/vjepa2}.}

\begin{itemize}
  \item We scale pretraining to 22M videos and a 1B parameter encoder.
  \item We show that the frozen encoder is a strong backbone for ``video question answering'' (see Fig.~\ref{fig:arch}).
\end{itemize}

% \section{Old section that was commented out}
\iffalse
This paragraph is disabled and must not be translated.
\fi

\section{Method}
\label{sec:method}
Given a video $x$, we mask a large fraction of the tokens and train the predictor to minimize
\begin{equation}
  \mathcal{L} = \lVert P_\phi(E_\theta(x_{\text{ctx}}), \Delta) - \text{sg}(\bar{E}(x)) \rVert_1 . \label{eq:loss}
\end{equation}
The target encoder $\bar{E}$ is an exponential moving average of $E_\theta$.
This prevents representation collapse without contrastive negatives.

\begin{figure}[h]
  \centering
  \fbox{\rule{0pt}{2cm}\rule{4cm}{0pt}}
  \caption{Overview of the architecture. The encoder processes visible patches; the predictor outputs the missing ones.}
  \label{fig:arch}
\end{figure}

\paragraph{Training details.} We use AdamW with a learning rate of $10^{-4}$ and a batch size of 3072 for 252k iterations.


\section{Experiments}
Table~\ref{tab:main} compares our approach with prior work, e.g. DINOv2 and SigLIP~2.
\begin{table}[t]
  \centering
  \caption{Results on Something-Something v2 and Diving-48. Higher is better.}
  \label{tab:main}
  \begin{tabular}{lcc}
    \toprule
    Method & SSv2 & Diving-48 \\
    \midrule
    DINOv2 & 50.7 & 79.3 \\
    \textbf{Ours} & \textbf{77.3} & \textbf{90.2} \\
    \bottomrule
  \end{tabular}
\end{table}

\begin{thebibliography}{9}
\bibitem{lecun2022} Y. LeCun. A path towards autonomous machine intelligence. 2022.
\bibitem{damen2022} D. Damen et al. Rescaling egocentric vision. IJCV, 2022.
\end{thebibliography}
\end{document}
